\section{Regressão Polinomial}
\label{sec:regressao}

Uma forma natural e eficaz de construir uma base é por meio de polinômios,
sendo que o conjunto \(\{\varphi_k\}\) pode ser formado por todas as funções
da forma
\begin{equation*}
    \varphi_k = \hat{e}_d \, c_{d,\vec{\alpha}} \, x^{\vec{\alpha}},
\end{equation*}
onde o vetor de expoentes $\vec{\alpha} = (\alpha_1, \alpha_2, \ldots, \alpha_n)
\in \mathbb{N}^n$ varia sobre todas as combinações cuja soma é
$|\vec{\alpha}| = \sum_{i=1}^n \alpha_i \le r$ --- com $r$ sendo o grau máximo
do polinômio ---, e $\hat{e}_d \in \mathbb{R}^m$ é o vetor unitário com valor
$1$ na componente de índice $d$ e zero nas demais.
Portanto, o polinômio pode ser escrito como

\begin{equation*}
    P_{\le r}(\vec{x}) =
        \sum_{d=1}^{m}\;
        \sum_{\substack{\vec{\alpha}\in\mathbb{N}^n\\|\vec{\alpha}|\le r}}
        c_{d,\vec{\alpha}}\;\hat{e}_d\,\vec{x}^{\,\vec{\alpha}},
        \qquad
        \vec{x}^{\,\vec{\alpha}}:=
        \prod_{i=1}^{n} x_i^{\alpha_i}.
\end{equation*}

onde \(\vec{x}^{\,\vec{\alpha}}\) denota o monômio associado ao multi-índice
\(\vec{\alpha}\), conforme \cite{burden_faires_2011}.

Expandindo os termos homogêneos $x^{\vec\alpha}$, obtemos

\begin{equation*}
    P_{\le r}(\vec{x}) =
    \sum_{d=1}^{m} \;
    \sum_{j=0}^{r} \;
    \sum_{\substack{\vec{\alpha}\in\mathbb{N}^n\\|\vec{\alpha}|=j}}
    c_{d,\vec{\alpha}}\;\hat{e}_d\,
    \prod_{i=1}^{n} x_i^{\alpha_i}.
\end{equation*}

Note que aqui assumimos a existência de algum mapeamento razoável na forma

\begin{equation*}
    (d,\vec\alpha) \to k.
\end{equation*}

O índice $k$ referencia cada função base de maneira unívoca e nos permite adotar
a modelagem matricial usual: montamos o sistema linear $G\vec{a} = \vec{b}$ em
que $G$ é a matriz de Gram associada às $\{\varphi_k\}$.

\begin{equation*}
    \begin{bmatrix}
    \langle \varphi_0, \varphi_0 \rangle &
    \langle \varphi_0, \varphi_1 \rangle &
    \cdots &
    \langle \varphi_0, \varphi_n \rangle \\
    \langle \varphi_1, \varphi_0 \rangle &
    \langle \varphi_1, \varphi_1 \rangle &
    \cdots &
    \langle \varphi_1, \varphi_n \rangle \\
    \vdots & \vdots & \ddots & \vdots \\
    \langle \varphi_n, \varphi_0 \rangle &
    \langle \varphi_n, \varphi_1 \rangle &
    \cdots &
    \langle \varphi_n, \varphi_n \rangle \\
    \end{bmatrix}
    \begin{bmatrix}
        a_0 \\ a_1 \\ \vdots \\ a_n
    \end{bmatrix}
    = \begin{bmatrix}
        \langle f,\varphi_0 \rangle \\
        \langle f,\varphi_1 \rangle \\
        \vdots \\
        \langle f,\varphi_n \rangle
    \end{bmatrix}
\end{equation*}

Essa formulação estabelece as relações que serão usadas nas seções seguintes
para analisar o condicionamento numérico e as estratégias de solução.
